%% toc: P15 171 171.5,171.6,171.7,171.8
\chapter{Uniform Convergence without a Countable Core, Bounded Losses and Lower Bounds}
\label{sec:p15-w02b-p15-065-beyond-core}

\Cref{sec:p15-vc-dimension-and-uniform-learning} proves a uniform deviation
bound for an event class of finite VC dimension under one measurability
hypothesis: a countable subfamily approximating every member pointwise
(\cref{def:p15-vcl-countable-core}). It notes that some classes of
dimension one, such as all singletons of an uncountable space, have no such
subfamily, and it treats only binary loss and only sufficient sample sizes.
This chapter supplies four further pieces.
\begin{enumerate}
\item A measurability regime that does not require a countable core. Its
      conclusions are stated in outer probability, so they remain meaningful
      when the supremum over the class or the output of a learner is not a
      measurable function of the sample. The regime asks only that the
      two-sample quantities produced by symmetrization be measurable for the
      completed product law. A worked class without a countable core
      satisfies it. Under the continuum hypothesis, a class of dimension one
      that violates it shows that the uniform conclusion can then fail.
\item The extension from binary loss to every loss with values in a bounded
      interval, through the class of level sets of the losses. The extension
      gives realizable and agnostic consequences and states how a positive
      noise floor rules out the realizable rate.
\item Lower bounds valid for every learner: a realizable bound of order
      \(d/\varepsilon\) and an agnostic bound of order \(d/\tau^2\), where
      \(d\) is the size of a shattered set. They are minimax statements, in
      which the sampling law may depend on the learner. They differ from
      the fixed-law countercases of
      \cref{thm:p15-vcl-examples}.
\item Conditions on the label noise that lie between the realizable and the
      agnostic regimes, how they imply one another, and the excess-risk
      rate each gives for empirical minimization, over a finite class and
      over a class of finite VC dimension.
\end{enumerate}

\emph{Inputs.} A probability space \((Z,\mathcal A,P)\); an event class or a
class of loss functions fixed before sampling; an independent sample of size
\(m\geq1\); for the lower bounds, a shattered finite set and an arbitrary
learner. \emph{Conclusions.} Explicit bounds on outer probabilities of
uniform deviation and of bad learner outputs, uniform over all members of the
class and over all laws satisfying the stated measurability condition; and,
for the lower bounds, a law under which a given learner fails with a stated
probability.

\section{Outer probability and measurability after symmetrization}
\label{sec:p15-w02b-p15-065-outer}

\begin{hbdefinition}{def:p15-w02b-p15-065-outer}{Outer probability, outer expectation and completion}
Let \((\Omega,\mathcal F,Q)\) be a probability space. For an arbitrary set
\(A\subseteq\Omega\) and an arbitrary function
\(f:\Omega\to[0,\infty)\), put
\begin{align*}
 Q^{*}(A)&=\inf\{Q(B):B\in\mathcal F,\ B\supseteq A\},\\
 E_Q^{*}f&=\inf\Bigl\{\int g\,dQ:\ g\ \mathcal F\text{-measurable},\ g\geq f\Bigr\}.
\end{align*}
These are the \emph{outer probability} and \emph{outer expectation}. The
\emph{completion} \(\overline{\mathcal F}\) consists of the sets \(A\) with
\(B_1\subseteq A\subseteq B_2\) for some \(B_1,B_2\in\mathcal F\) with
\(Q(B_2\setminus B_1)=0\), and \(\overline Q(A)=Q(B_2)\).
For a sample of size \(m\) from \(P\), \(P^{*}\) denotes the outer
probability of \(P^{\otimes m}\) on \((Z^m,\mathcal A^{\otimes m})\); for an
event that is measurable it is the ordinary probability.
\index{outer probability}\index{outer expectation}
\end{hbdefinition}

\begin{hbdefinition}{def:p15-w02b-p15-065-sym-measurable}{Two-sample quantities and measurability after symmetrization}
Let \(\mathcal C\subseteq\mathcal A\) be a nonempty event class fixed before
sampling, with trace growth \(\Pi_{\mathcal C}\) and empirical measure
\(P_m\) as in \cref{def:p15-vcl-event-class}. For
\(w=(z_1,\ldots,z_m,z'_1,\ldots,z'_m)\in Z^{2m}\) write \(P_m\) and
\(P'_m\) for the empirical measures of the first and second halves, and put
\begin{align*}
 D_m(w)&=\sup_{C\in\mathcal C}|P_m(C)-P'_m(C)|,\\
 K_m(w)&=\sup\bigl(\{0\}\cup\{P'_m(C):C\in\mathcal C,\ P_m(C)=0\}\bigr).
\end{align*}
The class is \emph{symmetrization-\allowbreak measurable for \(P\)} if, for every
\(m\geq1\), both \(D_m\) and \(K_m\) are measurable with respect to the
completion of \(\mathcal A^{\otimes2m}\) under \(P^{\otimes2m}\).
\index{symmetrization-measurable class}
\end{hbdefinition}

The condition concerns the two quantities that the ghost-sample argument
actually integrates. It does not require the one-sample discrepancy
\(\Delta_m\) to be measurable, and it does not require a countable core.
Measurability of \(\Delta_m\) alone is not enough:
\cref{thm:p15-w02b-p15-065-ch-failure} gives a class whose \(\Delta_m\)
is constant while its uniform deviation does not tend to zero.

\begin{hblemma}{thm:p15-w02b-p15-065-outer-fubini}{proved=annex:p15-w02b-p15-065-outer-fubini}{Outer expectation on a product space}
Let \((\Omega_1\times\Omega_2,\mathcal F_1\otimes\mathcal F_2,Q_1\otimes Q_2)\)
be a product of probability spaces and \(f:\Omega_1\times\Omega_2\to[0,1]\)
an arbitrary function. Then
\[
 E_{Q_1}^{*}\bigl[\omega_1\mapsto E_{Q_2}^{*}f(\omega_1,\cdot)\bigr]
 \leq E_{Q_1\otimes Q_2}^{*}f .
\]
Outer expectation is monotone and positively homogeneous; outer probability
is countably subadditive; and \(Q^{*}(A)=\overline Q(A)\) for every set
\(A\) in the completion. The reverse inequality can fail: under the
continuum hypothesis, item~4 of \cref{thm:p15-w02b-p15-065-ch-failure}
gives a set whose sections are all null and whose outer probability is one.
\end{hblemma}

\begin{hblemma}{thm:p15-w02b-p15-065-core-sufficient}{proved=annex:p15-w02b-p15-065-core-sufficient}{A countable core implies measurability after symmetrization}
If \(\mathcal C\) has a pointwise countable core \(\mathcal C_0\)
(\cref{def:p15-vcl-countable-core}), then \(D_m\) and \(K_m\) equal the
corresponding suprema over \(\mathcal C_0\). They are
\(\mathcal A^{\otimes2m}\)-measurable for every \(m\), and the class is
symmetrization-\allowbreak measurable for every law \(P\) on \((Z,\mathcal A)\).
\end{hblemma}

\begin{hbtheorem}{thm:p15-w02b-p15-065-outer-uniform}{proved=annex:p15-w02b-p15-065-outer-uniform}{Uniform deviation in outer probability}
Let \(\mathcal C\) be symmetrization-\allowbreak measurable for \(P\), let
\(Z_1,\ldots,Z_m\) be iid with law \(P\), \(m\geq1\), and let
\(\Delta_m=\sup_{C\in\mathcal C}|P_m(C)-P(C)|\).
\begin{enumerate}
\item For \(0<\varepsilon<1\),
      \(P^{*}(\Delta_m>\varepsilon)\leq
      \min\{1,\ 4\Pi_{\mathcal C}(2m)e^{-m\varepsilon^2/8}\}\).
\item If \(m\varepsilon\geq8\), then
      \(P^{*}\{\exists C\in\mathcal C:P_m(C)=0,\ P(C)>\varepsilon\}
      \leq2\Pi_{\mathcal C}(2m)2^{-m\varepsilon/2}\).
\end{enumerate}
The supremum ranges over every member of \(\mathcal C\), and the bounds do
not depend on \(P\). They therefore hold uniformly over every law for which
the class is symmetrization-\allowbreak measurable. When the class has a countable core,
the events are measurable and the statements reduce to
\cref{thm:p15-vcl-uniform,thm:p15-vcl-consistent}.
\end{hbtheorem}

The proof keeps the order of \cref{thm:p15-vcl-uniform}: a ghost sample,
then random swaps within a fixed pool, then a finite union over pooled
traces. The witness of a large deviation is now chosen sample by sample,
with no measurable selection. The first step needs no measurability at all,
because \cref{thm:p15-w02b-p15-065-outer-fubini} bounds an iterated outer
expectation by a joint one. The second step averages over the swaps, and
this is the only place where the completion-measurability of \(D_m\) and
\(K_m\) is used. The von Luxburg--Sch\"olkopf presentation states the
ghost-sample comparison as its Lemma~4 under the condition
\(m\varepsilon^2\geq2\), without proof
\cite[Section~5.2, p.~17]{LuxburgScholkopf2008Author}, and sets
measurability aside when it introduces the space of candidate functions
\cite[Section~2.2, p.~5]{LuxburgScholkopf2008Author}.

\begin{hbtheorem}{thm:p15-w02b-p15-065-outer-learning}{proved=annex:p15-w02b-p15-065-outer-learning}{Learning consequences without a countable core}
Let \(P\) be a law on \(X\times\{0,1\}\) and \(H\) a nonempty class of
measurable binary predictors fixed before sampling. Its loss-event class
\(\mathcal E_H=\{\{(x,y):h(x)\ne y\}:h\in H\}\) is assumed
symmetrization-\allowbreak measurable for \(P\); no countable core is assumed. Use
\(R\), \(\widehat R_m\) and \(R_H^{*}\) as in
\cref{def:p15-vcl-classification}, and let \(\Pi_H\) be the growth of the
positive prediction-event class.
\begin{enumerate}
\item \(\Pi_{\mathcal E_H}=\Pi_H\). For every \(\eta\geq0\) and every
      sample, an \(\eta\)-approximate empirical minimizer exists, and an
      exact one exists. Let \(\widehat h\) be any map assigning to each
      sample such a minimizer, measurable or not. Then for
      \(0<\varepsilon<1\),
      \[
       P^{*}\{R(\widehat h)>R_H^{*}+2\varepsilon+\eta\}
       \leq\min\{1,\ 4\Pi_H(2m)e^{-m\varepsilon^2/8}\}.
      \]
\item If some \(h^{*}\in H\) has \(R(h^{*})=0\), let \(\widehat h\) be any map
      returning a member of \(H\) with no training error whenever one
      exists. Then for \(m\varepsilon\geq8\),
      \(P^{*}\{R(\widehat h)>\varepsilon\}\leq2\Pi_H(2m)2^{-m\varepsilon/2}\).
\item The sufficient sample sizes of \cref{thm:p15-vcl-samples}, for
      \(\Pi_H\) of finite dimension \(d\geq1\), remain sufficient: the first
      formula for item~1 and the second for item~2.
\end{enumerate}
The validation statement of \cref{thm:p15-vcl-erm} is not extended to
non-measurable learners: \cref{thm:p15-w02b-p15-065-ch-failure} shows that
it can then fail.
\end{hbtheorem}

\begin{hbtheorem}{thm:p15-w02b-p15-065-singletons}{proved=annex:p15-w02b-p15-065-singletons}{Singletons: dimension one, no countable core, measurable after symmetrization}
Let \(Z=[0,1]\) with its Borel sets, and let
\(\mathcal C=\{\{x\}:x\in[0,1]\}\cup\{\varnothing\}\).
\begin{enumerate}
\item \(\mathcal C\) has VC dimension one, \(\Pi_{\mathcal C}(t)=t+1\) for
      \(t\geq1\), and no pointwise countable core.
\item For every Borel law \(P\) on \([0,1]\) and every \(m\), the functions
      \(D_m\), \(K_m\) and \(\Delta_m\) are Borel measurable. The class is
      symmetrization-\allowbreak measurable for every such \(P\), and
      \cref{thm:p15-w02b-p15-065-outer-uniform} applies with
      \(\Pi_{\mathcal C}(2m)=2m+1\).
\item Under the uniform law, \(\Delta_m=1/m\) almost surely. Under a point
      mass, \(\Delta_m=0\) almost surely.
\end{enumerate}
\end{hbtheorem}

\begin{hbtheorem}{thm:p15-w02b-p15-065-ch-failure}{proved=annex:p15-w02b-p15-065-ch-failure}{Under the continuum hypothesis: dimension one without uniform convergence}
Assume the continuum hypothesis. There is a well-ordering \(\prec\) of
\([0,1]\) in which every set \(I_x=\{y:y\preceq x\}\) is countable. Let
\(P\) be the uniform law on \([0,1]\) and
\(\mathcal C=\{I_x:x\in[0,1]\}\).
\begin{enumerate}
\item Each \(I_x\) is a Borel set of probability zero. \(\mathcal C\) has
      VC dimension one and \(\Pi_{\mathcal C}(t)\leq t+1\).
\item For every \(m\geq1\) and every sample, \(\Delta_m=1\).
\item Consequently \(\mathcal C\) has no pointwise countable core, and it is
      not symmetrization-\allowbreak measurable for \(P\). In fact, for every \(m\)
      with \(4(2m+1)e^{-m/32}<1\) the two-sample discrepancy \(D_m\) is
      not measurable for the completed product law, although \(\Delta_m\)
      is constant.
\item Take labels \(Y\equiv0\), training sample size one and the learner
      that returns \(\widehat h_y(x)=\mathbf1\{x\prec y\}\) on training
      input \(y\). For every \(v\geq1\) and every \(0\leq u<1\), each fixed
      training input gives validation probability zero to
      \(\{\widehat R_v(\widehat h_y)>u\}\), while
      \(R(\widehat h_y)=0\). The outer probability of that event under the
      joint law of training and validation samples is one.
\end{enumerate}
\end{hbtheorem}

\begin{hbwarning}{def:p15-w02b-p15-065-measurability-warning}{What the measurability condition does and does not buy}
Finite VC dimension is a combinatorial property of traces. It controls
deviations only together with a measurability condition linking the class
to the sampling law. \Cref{thm:p15-w02b-p15-065-ch-failure} uses the
continuum hypothesis, so it shows that the condition cannot be dropped from
a theorem provable in ordinary set theory. It does not show that every
class without a countable core fails;
\cref{thm:p15-w02b-p15-065-singletons} is a counterexample to that
reading. Outer probability bounds the size of a bad event even when the
event is not measurable. It does not turn a section-by-section bound into
a joint one, which is why the validation bound needs a jointly measurable
learner.
\end{hbwarning}

\section{Losses with values in a bounded interval}
\label{sec:p15-w02b-p15-065-real}

\begin{hbdefinition}{def:p15-w02b-p15-065-level-class}{Loss class, level-set class and its capacity}
Let \(\mathcal G\) be a nonempty class of measurable functions
\(g:Z\to[0,1]\) fixed before sampling. In learning, \(Z=X\times\mathcal Y\)
and \(g_h(x,y)=\ell(h(x),y)\) for a loss \(\ell\) with values in \([0,1]\)
and a predictor \(h\) in a class \(H\). Write
\(P_mg=m^{-1}\sum_ig(Z_i)\) and \(Pg=\int g\,dP\); for learning,
\(\widehat R_m(h)=P_mg_h\), \(R(h)=Pg_h\) and
\(R_H^{*}=\inf_{h\in H}R(h)\). The \emph{level-set class} is
\[
 \mathcal L_{\mathcal G}=\{\{z:g(z)>s\}:g\in\mathcal G,\ 0<s<1\},
\]
with growth \(\Pi_{\mathcal L_{\mathcal G}}\) and VC dimension
\(d_{\mathcal G}\) in the sense of \cref{def:p15-vcl-event-class}.
An \(\eta\)-approximate empirical minimizer is a member \(\widehat h\) with
\(\widehat R_m(\widehat h)\leq\inf_{h\in H}\widehat R_m(h)+\eta\).
\end{hbdefinition}

\begin{hbtheorem}{thm:p15-w02b-p15-065-real-uniform}{proved=annex:p15-w02b-p15-065-real-uniform}{Uniform deviation for bounded losses}
Let \(Z_1,\ldots,Z_m\) be iid with law \(P\).
\begin{enumerate}
\item On every sample,
      \(\sup_{g\in\mathcal G}|P_mg-Pg|\leq
      \sup_{L\in\mathcal L_{\mathcal G}}|P_m(L)-P(L)|\).
\item If \(\mathcal L_{\mathcal G}\) is symmetrization-\allowbreak measurable for \(P\),
      then for \(0<\varepsilon<1\),
      \[
       P^{*}\Bigl\{\sup_{g\in\mathcal G}|P_mg-Pg|>\varepsilon\Bigr\}
       \leq\min\{1,\ 4\Pi_{\mathcal L_{\mathcal G}}(2m)e^{-m\varepsilon^2/8}\}.
      \]
\item For functions with values in \([a,b]\), \(a<b\), apply items~1
      and~2 to \((g-a)/(b-a)\). The level sets become
      \(\{g>r\}\) with \(a<r<b\). For \(0<\varepsilon<b-a\) the bound
      becomes \(4\Pi(2m)e^{-m\varepsilon^2/(8(b-a)^2)}\).
\item If \(m\varepsilon\geq8\) and \(\mathcal L_{\mathcal G}\) is
      symmetrization-\allowbreak measurable for \(P\), then
      \(P^{*}\{\exists g\in\mathcal G:P_mg=0,\ Pg>\varepsilon\}
      \leq2\Pi_{\mathcal L_{\mathcal G}}(2m)2^{-m\varepsilon/2}\).
\item If \(1\leq d_{\mathcal G}<\infty\), the sample sizes of
      \cref{thm:p15-vcl-samples} with \(d=d_{\mathcal G}\) suffice with the
      same \(\varepsilon\) and \(\delta\): the first formula for item~2 and
      the second for item~4.
\end{enumerate}
\end{hbtheorem}

\begin{hbtheorem}{thm:p15-w02b-p15-065-real-learning}{proved=annex:p15-w02b-p15-065-real-learning}{Empirical minimization for bounded losses}
In the learning reading of \cref{def:p15-w02b-p15-065-level-class},
suppose \(\mathcal L_{\mathcal G}\) is symmetrization-\allowbreak measurable for \(P\).
\begin{enumerate}
\item For every \(\eta>0\) and every sample an \(\eta\)-approximate
      empirical minimizer exists. An exact minimizer need not exist.
\item For any map \(\widehat h\) returning an \(\eta\)-approximate
      empirical minimizer, \(\eta\geq0\), and \(0<\varepsilon<1\),
      \[
       P^{*}\{R(\widehat h)>R_H^{*}+2\varepsilon+\eta\}
       \leq\min\{1,\ 4\Pi_{\mathcal L_{\mathcal G}}(2m)e^{-m\varepsilon^2/8}\}.
      \]
\item If some \(h^{*}\in H\) has \(R(h^{*})=0\), then almost surely
      \(\widehat R_m(h^{*})=0\). For any map returning a member with zero
      empirical risk whenever one exists and \(m\varepsilon\geq8\),
      \(P^{*}\{R(\widehat h)>\varepsilon\}
      \leq2\Pi_{\mathcal L_{\mathcal G}}(2m)2^{-m\varepsilon/2}\).
\end{enumerate}
\end{hbtheorem}

\begin{hbtheorem}{thm:p15-w02b-p15-065-noise}{proved=annex:p15-w02b-p15-065-noise}{Noise floor and the two regimes for bounded losses}
Consider a prediction problem with loss in \([0,1]\), minimum achievable
risk \(R^{*}\) and class \(H\), as in
\cref{def:p15-w02b-p15-064-risk} with \(\mathcal F=H\).
\begin{enumerate}
\item If some member of \(H\) has risk zero, then \(R^{*}=0\). Hence a
      positive noise floor \(R^{*}>0\) excludes the realizable regime for
      every class, and of the two bounds in
      \cref{thm:p15-w02b-p15-065-real-learning} only item~2 applies.
\item \(R^{*}=0\) does not imply that \(H\) contains a member of risk zero.
\item On the event of item~2 of \cref{thm:p15-w02b-p15-065-real-learning},
      \(R(\widehat h)-R^{*}\leq(R_H^{*}-R^{*})+2\varepsilon+\eta\). The first
      term does not depend on the sample or on \(m\).
\end{enumerate}
\end{hbtheorem}

In the realizable regime, item~3 of
\cref{thm:p15-w02b-p15-065-real-learning} gives an exponent linear in
\(m\varepsilon\), so the sufficient sample size grows like \(1/\varepsilon\)
up to logarithmic factors. Under noise only the exponent
\(m\varepsilon^2/8\) is available, so the size grows like \(1/\varepsilon^2\).
\Cref{thm:p15-w02b-p15-065-lower-realizable,thm:p15-w02b-p15-065-lower-agnostic}
show that for binary loss, which has values in \([0,1]\), neither order
can be improved beyond logarithmic factors uniformly over laws. Conditions
on the noise that lie between these two regimes, and the rates they give,
are the subject of \cref{sec:p15-w02b-p15-065-noise-conditions}.

\begin{hbexample}{def:p15-w02b-p15-065-absolute-example}{Absolute loss with constant predictions}
Let \(Y\) be uniform on \([0,1]\) and independent of an input \(X\). Let
\(H\) be the constant predictors \(\theta\in[0,1]\) and take absolute loss,
so \(g_\theta(x,y)=|y-\theta|\in[0,1]\).
\Cref{thm:p15-w02b-p15-065-absolute-example} proves the following.
\(\mathcal L_{\mathcal G}\) has VC dimension two and a countable core, and
\(\Pi_{\mathcal L_{\mathcal G}}(t)\leq(t^2+t+2)/2\). The risk is
\(R(\theta)=(\theta^2+(1-\theta)^2)/2\), so \(R_H^{*}=R^{*}=1/4\),
attained at \(\theta=1/2\), and the problem is not realizable. With
\(\varepsilon=0.1\) and \(\delta=0.05\), the first formula of
\cref{thm:p15-vcl-samples} with \(d=2\) gives
\[
 m\geq\Bigl\lceil 6400\bigl(2\ln6400+\ln160\bigr)\Bigr\rceil=144661 .
\]
For such \(m\) every \(\eta\)-approximate empirical minimizer satisfies
\(P^{*}\{R(\widehat\theta)>0.45+\eta\}\leq0.05\). The constants are
sufficient, not sharp.
\end{hbexample}

\begin{hbtheorem}{thm:p15-w02b-p15-065-absolute-example}{proved=annex:p15-w02b-p15-065-absolute-example}{The calculations of the absolute-loss example}
Every assertion of \cref{def:p15-w02b-p15-065-absolute-example} holds.
\end{hbtheorem}

\section{Lower bounds for every learner}
\label{sec:p15-w02b-p15-065-lower}

\begin{hbdefinition}{def:p15-w02b-p15-065-learner}{Learners and finitely supported laws}
Let \(X\) be a measurable space whose singletons are measurable, and \(H\)
a class of binary predictors. A \emph{learner} of sample size \(m\) is any
map \(A\) from \((X\times\{0,1\})^m\) to the set of all functions
\(X\to\{0,1\}\). Its output need not lie in \(H\) and need not be
measurable. A randomized learner also takes an independent internal draw
\(U\), and \(A(s,u)(x)\) is required to be measurable in \(u\) for each
\(s\) and \(x\). For a law \(P\) supported on finitely many points
\((x,y)\), \(R_P(f)=\sum_{(x,y)}P\{(x,y)\}\mathbf1\{f(x)\ne y\}\) is
defined for every function \(f\). The sample takes finitely many values
almost surely, so every event about \(R_P(A(S))\) is measurable.
\end{hbdefinition}

\begin{hbtheorem}{thm:p15-w02b-p15-065-lower-realizable}{proved=annex:p15-w02b-p15-065-lower-realizable}{Realizable lower bound}
Let \(d\geq2\) and let \(H\) shatter \(d\) distinct points. Let
\(0<\varepsilon\leq1/8\) and let \(m\) be an integer with
\(1\leq m\leq(d-1)/(16\varepsilon)\). For every learner \(A\), possibly
randomized, there is a law \(P\) supported on those \(d\) points and a
member \(h^{*}\in H\) with \(R_P(h^{*})=0\) such that
\[
 P^{\otimes m}\{R_P(A(S))\geq\varepsilon\}\geq\tfrac18 .
\]
Hence no learner achieves risk below \(\varepsilon\) with probability more
than \(7/8\) for every realizable law once \(m\leq(d-1)/(16\varepsilon)\).
\end{hbtheorem}

\begin{hbtheorem}{thm:p15-w02b-p15-065-lower-agnostic}{proved=annex:p15-w02b-p15-065-lower-agnostic}{Agnostic lower bound}
Let \(d\geq1\) and let \(H\) shatter \(d\) distinct points. Let
\(0<\tau\leq1/16\) and let \(m\) be an integer with
\(1\leq m\leq d/(512\tau^2)\). For every learner \(A\), possibly
randomized, there is a law \(P\) supported on those \(d\) points with both
labels such that
\[
 P^{\otimes m}\{R_P(A(S))-R_H^{*}\geq\tau\}\geq\tfrac1{16},
\]
where \(R_H^{*}=\inf_{h\in H}R_P(h)\) equals the minimum achievable risk of
\(P\) and is attained in \(H\).
\end{hbtheorem}

\begin{hbcorollary}{thm:p15-w02b-p15-065-lower-infinite}{proved=annex:p15-w02b-p15-065-lower-infinite}{Infinite dimension excludes distribution-free sample sizes}
If \(H\) has infinite VC dimension, then for every \(m\geq1\), every
\(0<\varepsilon\leq1/8\) and every learner, some realizable finitely
supported law gives \(P^{\otimes m}\{R_P(A(S))\geq\varepsilon\}\geq1/8\).
No finite sample size achieves accuracy \(\varepsilon\) with confidence
\(\delta<1/8\) for every realizable law, and hence none does so for every
law.
\end{hbcorollary}

\begin{hbwarning}{def:p15-w02b-p15-065-minimax-warning}{Minimax necessity is not fixed-law necessity}
In \cref{thm:p15-w02b-p15-065-lower-realizable,thm:p15-w02b-p15-065-lower-agnostic,thm:p15-w02b-p15-065-lower-infinite}
the law is chosen after the learner. For each fixed law in those
constructions, the learner that ignores the data and outputs that law's
best predictor has excess risk zero. The bounds therefore state what no
single learner can guarantee for all laws together. They do not say that
learning fails for a particular law. This is a different kind of statement
from \cref{thm:p15-vcl-examples}, item~2, which compares two fixed laws
for one countable class and concerns uniform deviation, not learners.
Combined with \cref{thm:p15-vcl-samples}, which needs its measurability
hypothesis and a class of VC dimension \(d\), the realizable bound shows
that sample size of order \(d/\varepsilon\) is necessary and, up to the
logarithmic factor in the sufficient formula, enough. The agnostic bound
shows the same for order \(d/\tau^2\). The constants on both sides are not
optimized.
\end{hbwarning}

\section{Noise conditions between the realizable and agnostic regimes}
\label{sec:p15-w02b-p15-065-noise-conditions}

The realizable rate of order \(1/\varepsilon\) and the agnostic rate of
order \(1/\varepsilon^2\) are the two ends of a scale. Which end applies is
governed by how often the label is close to a fair coin given the input.
The von Luxburg--Sch\"olkopf presentation describes learning as harder the
closer the conditional label probability lies to one half
\cite[Section~2.2, p.~3]{LuxburgScholkopf2008Author}, and remarks that
assumptions on the label noise improve the rates for the estimation error
\cite[Section~7.4, p.~29]{LuxburgScholkopf2008Author}. This section states
three such assumptions, proves how they are related, and proves the rate
each gives for empirical minimization over a finite class and over a class
of finite VC dimension.

\emph{Inputs.} A law \(P\) on \(\mathcal X\times\{0,1\}\), a class of
measurable binary predictors fixed before sampling, either finite or of
finite VC dimension with the measurability hypothesis of
\cref{def:p15-w02b-p15-065-ternary}, an iid sample, and one of the
conditions of \cref{def:p15-w02b-p15-065-noise-conditions}.
\emph{Conclusion.} A bound on the excess risk of every empirical
minimizer, holding outside a set of outer probability at most
\(\delta\), whose order in \(m\) moves from \(m^{-1/2}\) to \(m^{-1}\), up
to the logarithmic factor of the growth for an infinite class, as the
condition strengthens.

\begin{hbdefinition}{def:p15-w02b-p15-065-noise-conditions}{Conditions on the label noise}
Let \(\mathcal X\) be a measurable space and \(P\) the law of a pair
\((X,Y)\) with values in \(\mathcal X\times\{0,1\}\). Let \(\eta\) be a
measurable function on \(\mathcal X\) such that \(\eta(X)\) is a version
of \(P(Y=1\mid X)\). Let \(f_{\mathrm{Bayes}}\) and the minimum achievable
risk \(R^{*}\) be as in
\cref{thm:p15-w02b-p15-064-classification-bayes}, with labels
\(\{0,1\}\). For a measurable \(f:\mathcal X\to\{0,1\}\) write
\(R(f)=P\{f(X)\ne Y\}\).
\begin{enumerate}
\item \(P\) satisfies the \emph{bounded-noise condition with margin
      \(\gamma\)}, \(0<\gamma\leq1\), if
      \(|2\eta(X)-1|\geq\gamma\) almost surely.
\item \(P\) satisfies the \emph{margin condition with exponent \(a>0\) and
      constant \(C>0\)} if
      \(P\{|2\eta(X)-1|\leq t\}\leq Ct^{a}\) for every \(0<t\leq1\).
      Taking \(t=1\) shows that then \(C\geq1\).
\item Let \(\mathcal F\) be a set of measurable predictors and
      \(f^{\circ}\in\mathcal F\) with
      \(R(f^{\circ})=\min_{f\in\mathcal F}R(f)\). \(P\) satisfies the
      \emph{variance condition on \(\mathcal F\) with exponent
      \(\kappa\in[0,1]\) and constant \(B>0\), relative to \(f^{\circ}\)},
      if for every \(f\in\mathcal F\)
      \[
       P\{f(X)\ne f^{\circ}(X)\}\leq B\bigl(R(f)-R(f^{\circ})\bigr)^{\kappa},
      \]
      with the convention \(0^0=1\).
\end{enumerate}
The problem is \emph{realizable by \(f^{\circ}\)} if \(Y=f^{\circ}(X)\)
almost surely.
\index{bounded-noise condition}\index{margin condition}\index{variance condition}
\end{hbdefinition}

The variance condition is the one the proof uses. It bounds how often a
predictor disagrees with the reference member by a power of how much worse
it is. With \(\kappa=0\) and \(B=1\) it says nothing; with \(\kappa=1\) a
predictor can disagree with the reference only on a set whose probability
is proportional to its excess risk. The first two conditions are stated
on \(\eta\) and translate into it.

\begin{hblemma}{thm:p15-w02b-p15-065-noise-relations}{proved=annex:p15-w02b-p15-065-noise-relations}{How the three conditions are related}
Assume the setting of \cref{def:p15-w02b-p15-065-noise-conditions}, and
let \(\mathcal F_{\mathrm{all}}\) be the set of all measurable predictors.
\begin{enumerate}
\item If the problem is realizable by a measurable \(f^{\circ}\), the
      bounded-noise condition holds with margin \(\gamma=1\).
\item The bounded-noise condition with margin \(\gamma\) implies the
      margin condition with every exponent \(a>0\) and constant
      \(C=\gamma^{-a}\).
\item Under the bounded-noise or the margin condition,
      \(P\{\eta(X)=1/2\}=0\), and every Bayes predictor \(f^{*}\) agrees
      with \(f_{\mathrm{Bayes}}\) almost surely.
\item The bounded-noise condition with margin \(\gamma\) implies the
      variance condition on \(\mathcal F_{\mathrm{all}}\) with \(\kappa=1\)
      and \(B=1/\gamma\), relative to every Bayes predictor.
\item The margin condition with exponent \(a\) and constant \(C\) implies
      the variance condition on \(\mathcal F_{\mathrm{all}}\) with
      \(\kappa=a/(1+a)\) and \(B=2C^{1/(1+a)}\), relative to every Bayes
      predictor.
\item Every law satisfies the variance condition on every \(\mathcal F\),
      relative to every minimizing member, with \(\kappa=0\) and \(B=1\).
\item If \(P\{\eta(X)=1/2\}>0\), no variance condition on
      \(\mathcal F_{\mathrm{all}}\) with \(\kappa>0\) holds relative to any
      Bayes predictor.
\item A variance condition on \(\mathcal F_{\mathrm{all}}\) relative to a
      Bayes predictor \(f^{*}\) restricts to every class \(\mathcal F\)
      containing \(f^{*}\), with the same \(\kappa\) and \(B\).
\end{enumerate}
\end{hblemma}

\begin{hbtheorem}{thm:p15-w02b-p15-065-noise-rate}{proved=annex:p15-w02b-p15-065-noise-rate}{Excess risk of empirical minimization under a variance condition}
Let \(H\) be a finite class of \(r\geq1\) measurable predictors
\(\mathcal X\to\{0,1\}\) fixed before sampling, let \((X_1,Y_1),\ldots,(X_m,Y_m)\)
be iid with law \(P\), \(m\geq1\), and write
\(\widehat R_m(h)=m^{-1}\sum_i\mathbf1\{h(X_i)\ne Y_i\}\) and
\(R_H^{*}=\min_{h\in H}R(h)\). Suppose \(P\) satisfies the variance
condition on \(H\) with exponent \(\kappa\) and constant \(B\), relative to
\(h^{\circ}\in H\). For \(0<\delta<1\) put \(L=\ln(r/\delta)\). There is an
event of probability at least \(1-\delta\) on which every \(h\in H\)
satisfies
\[
 R(h)-R_H^{*}\leq\widehat R_m(h)-\widehat R_m(h^{\circ})
 +\sqrt{\frac{2B\,(R(h)-R_H^{*})^{\kappa}L}{m}}+\frac{4L}{3m}.
\]
On that event:
\begin{enumerate}
\item every \(h\in H\) with \(\widehat R_m(h)\leq\widehat R_m(h^{\circ})\),
      in particular every exact empirical minimizer, satisfies
      \[
       R(h)-R_H^{*}\leq\max\Bigl\{\Bigl(\frac{8BL}{m}\Bigr)^{1/(2-\kappa)},\
       \frac{8L}{3m}\Bigr\};
      \]
\item every \(\vartheta\)-approximate empirical minimizer,
      \(\vartheta\geq0\), satisfies
      \[
       R(h)-R_H^{*}\leq\max\Bigl\{3\vartheta,\
       \Bigl(\frac{18BL}{m}\Bigr)^{1/(2-\kappa)},\ \frac{4L}{m}\Bigr\}.
      \]
\end{enumerate}
The selection among minimizers may be arbitrary; no measurability of the
selected member is needed. If \(h^{\circ}\) is a Bayes predictor, then
\(R_H^{*}=R^{*}\) and the bounds are bounds on the excess risk over the
minimum achievable risk.
\end{hbtheorem}

\begin{hbcorollary}{thm:p15-w02b-p15-065-noise-sizes}{proved=annex:p15-w02b-p15-065-noise-sizes}{Sufficient sample sizes across the noise scale}
In the setting of \cref{thm:p15-w02b-p15-065-noise-rate}, let
\(0<\varepsilon\leq1\) and \(L=\ln(r/\delta)\). Exact empirical
minimization has \(R(\widehat h)-R_H^{*}\leq\varepsilon\) with probability
at least \(1-\delta\) whenever
\(m\geq\max\{8BL/\varepsilon^{2-\kappa},\,8L/(3\varepsilon)\}\). In
particular, when \(H\) contains a Bayes predictor, so that
\(R_H^{*}=R^{*}\), each of the following sizes suffices.
\begin{enumerate}
\item Realizable by a member of \(H\): \(m\geq8L/\varepsilon\).
\item Bounded noise with margin \(\gamma\):
      \(m\geq8L/(\gamma\varepsilon)\).
\item Margin condition with exponent \(a\) and constant \(C\):
      \[
       m\geq16C^{1/(1+a)}L\,\varepsilon^{-(2+a)/(1+a)}.
      \]
\item No condition on the noise: \(m\geq8L/\varepsilon^2\).
\end{enumerate}
The exponent \((2+a)/(1+a)\) of item~3 decreases from \(2\) towards \(1\)
as \(a\) increases from \(0\) to \(\infty\).
\end{hbcorollary}

\begin{hbcorollary}{thm:p15-w02b-p15-065-noise-lower}{proved=annex:p15-w02b-p15-065-noise-lower}{The margin enters the sample size}
Let \(d\geq1\), let \(H\), a class of measurable binary predictors,
shatter \(d\) distinct points of a space with measurable singletons, and
let \(0<\gamma\leq1\). For every learner \(A\)
of sample size \(m\) with \(1\leq m\leq d/(2\gamma^2)\), possibly
randomized, there is a law \(P\) supported on those \(d\) points that
satisfies the bounded-noise condition with margin \(\gamma\), for which
some member of \(H\) is a Bayes predictor, and with
\[
 P^{\otimes m}\{R_P(A(S))-R^{*}\geq\gamma/16\}\geq\tfrac1{16}.
\]
With \(\varepsilon=\gamma/16\) the size condition reads
\(m\leq d/(32\gamma\varepsilon)\).
\end{hbcorollary}

For a finite class shattering \(d\) points, \(r\geq2^d\), so
\(L\geq d\ln2\). Item~2 of \cref{thm:p15-w02b-p15-065-noise-sizes} and
\cref{thm:p15-w02b-p15-065-noise-lower} then show that, when the excess
risk asked for is comparable to the margin, the factor
\(1/(\gamma\varepsilon)\) in the sufficient size cannot be removed by any
learner. The lower bound does not cover small \(\varepsilon\) at fixed
\(\gamma\).

\begin{hbexample}{def:p15-w02b-p15-065-noise-example}{Thresholds with bounded noise}
Let \(X\) be uniform on \([0,1]\) and \(\eta(x)=1/10\) for \(x<3/10\),
\(\eta(x)=9/10\) for \(x\geq3/10\). Let \(H=\{h_0,\ldots,h_{100}\}\) with
\(h_j(x)=\mathbf1\{x\geq j/100\}\), so \(r=101\).
\Cref{thm:p15-w02b-p15-065-noise-example} proves the following. The member
\(h_{30}\) is a Bayes predictor and \(R^{*}=1/10\). The bounded-noise
condition holds with \(\gamma=4/5\), and the variance condition on \(H\)
holds relative to \(h_{30}\) with \(\kappa=1\) and \(B=5/4\), with equality
for every member. With \(\varepsilon=1/100\) and \(\delta=1/20\),
\(L=\ln2020\), and item~2 of \cref{thm:p15-w02b-p15-065-noise-sizes}
asks for \(m\geq7611\). The bound that ignores the noise structure,
\cref{thm:p15-conc-finite-learning}, asks for \(m\geq166080\) for the same
\(\varepsilon\) and \(\delta\); item~4 of
\cref{thm:p15-w02b-p15-065-noise-sizes} asks for \(m\geq608869\). Every
exact empirical minimizer over a sample of size \(7611\) has risk at most
\(0.11\) with probability at least \(0.95\).
\end{hbexample}

\begin{hbtheorem}{thm:p15-w02b-p15-065-noise-example}{proved=annex:p15-w02b-p15-065-noise-example}{The calculations of the threshold example}
Every assertion of \cref{def:p15-w02b-p15-065-noise-example} holds.
\end{hbtheorem}

The finite-class proof applies Bernstein's inequality to each of the \(r\)
loss differences and takes a union bound. A class of finite VC dimension
may have infinitely many members, so the union is replaced by a ghost
sample, random swaps and a count of the patterns the loss differences leave
on the pooled sample, as in \cref{thm:p15-w02b-p15-065-outer-uniform}. The
variance enters through the number of pooled points at which a member
disagrees with the reference member: the swap bound is normalized by that
count instead of by \(m\). The measurability hypothesis is of the same kind
as in \cref{def:p15-w02b-p15-065-sym-measurable}: one function of the
pooled sample must be measurable for the completed product law.

\begin{hbdefinition}{def:p15-w02b-p15-065-ternary}{Ternary classes and a normalized two-sample discrepancy}
Let \(\mathcal T\) be a nonempty class of measurable functions
\(g:Z\to\{-1,0,1\}\) fixed before sampling. For \(t\geq1\) let
\(\Pi_{\mathcal T}(t)\) be the largest number of distinct vectors
\((g(z_1),\ldots,g(z_t))\), \(g\in\mathcal T\), on any ordered \(t\)-tuple.
For \(w=(z_1,\ldots,z_m,z'_1,\ldots,z'_m)\in Z^{2m}\), with \(P_m\) and
\(P'_m\) the empirical measures of the two halves as in
\cref{def:p15-w02b-p15-065-sym-measurable}, put
\begin{gather*}
 n_g(w)=m\bigl(P_m|g|+P'_m|g|\bigr),\\
 J_m(w)=\max\Bigl\{0,\ \sup_{g\in\mathcal T,\ n_g(w)>0}
 \frac{m(P'_mg-P_mg)}{\sqrt{n_g(w)}}\Bigr\},
\end{gather*}
with the supremum of the empty set read as \(0\). Thus \(n_g(w)\) is the
number of the \(2m\) coordinates of \(w\) at which \(g\) is not zero. The
class is \emph{self-normalized symmetrization-\allowbreak measurable for
\(P\) at sample size \(m\)} if \(J_m\) is measurable with respect to the
completion of \(\mathcal A^{\otimes2m}\) under \(P^{\otimes2m}\).

In the setting of \cref{def:p15-w02b-p15-065-noise-conditions}, with
\(Z=\mathcal X\times\{0,1\}\), a class \(H\) of measurable predictors and
a reference member \(h^{\circ}\in H\), put for \(h\in H\)
\[
 \xi_h(x,y)=\mathbf1\{h(x)\ne y\}-\mathbf1\{h^{\circ}(x)\ne y\},\qquad
 \chi_h(x,y)=\mathbf1\{h(x)\ne h^{\circ}(x)\},
\]
and \(\mathcal T_H=\{\xi_h:h\in H\}\cup\{-\chi_h:h\in H\}\). Write
\(\Pi_H\) for the growth of the positive prediction-event class
\(\{\{x:h(x)=1\}:h\in H\}\), as in \cref{def:p15-vcl-classification}.
\index{self-normalized symmetrization-measurable class}
\end{hbdefinition}

\begin{hblemma}{thm:p15-w02b-p15-065-ternary-core}{proved=annex:p15-w02b-p15-065-ternary-core}{Patterns and measurability of the ternary class of a predictor class}
In the second part of \cref{def:p15-w02b-p15-065-ternary}:
\begin{enumerate}
\item \(\xi_h^2=|\xi_h|=\chi_h\), so \(P\chi_h=P\{h(X)\ne h^{\circ}(X)\}\)
      and \(P\xi_h=R(h)-R(h^{\circ})\).
\item \(\Pi_{\mathcal T_H}(t)\leq2\Pi_H(t)\) for every \(t\geq1\).
\item If the positive prediction-event class of \(H\) has a pointwise
      countable core (\cref{def:p15-vcl-countable-core}), then for every
      \(m\) the function \(J_m\) of \(\mathcal T_H\) equals the
      corresponding expression with \(H\) replaced by the core, it is
      measurable for \(\mathcal A^{\otimes2m}\), and \(\mathcal T_H\) is
      self-normalized symmetrization-\allowbreak measurable for every law
      \(P\) and every \(m\).
\end{enumerate}
\end{hblemma}

\begin{hblemma}{thm:p15-w02b-p15-065-relative}{proved=annex:p15-w02b-p15-065-relative}{Deviation normalized by the frequency of nonzero values}
Let \(\mathcal T\) be as in \cref{def:p15-w02b-p15-065-ternary}, let
\(Z_1,\ldots,Z_m\) be iid with law \(P\), \(m\geq1\), and suppose
\(\mathcal T\) is self-normalized symmetrization-\allowbreak measurable
for \(P\) at sample size \(m\). For every \(L>0\),
\begin{multline*}
 P^{*}\Bigl\{\exists g\in\mathcal T:\ Pg-P_mg>
 \sqrt{\frac{80\max\{P|g|,\,P_m|g|\}\,L}{m}}\Bigr\}\\
 \leq\min\bigl\{1,\ 2\Pi_{\mathcal T}(2m)e^{-L}\bigr\}.
\end{multline*}
\end{hblemma}

The normalization is by the larger of the true and the empirical
frequency of nonzero values, not by \(1\). For a member that is rarely
nonzero the permitted deviation is correspondingly small; this is what
turns a variance condition into a faster rate.

\begin{hbtheorem}{thm:p15-w02b-p15-065-vc-noise-rate}{proved=annex:p15-w02b-p15-065-vc-noise-rate}{Excess risk of empirical minimization over a class of finite VC dimension under a variance condition}
Let \(P\) be a law on \(\mathcal X\times\{0,1\}\) and \(H\) a nonempty
class of measurable predictors \(\mathcal X\to\{0,1\}\) fixed before
sampling, possibly infinite. Suppose \(h^{\circ}\in H\) has
\(R(h^{\circ})=R_H^{*}=\inf_{h\in H}R(h)\), and \(P\) satisfies the variance
condition on \(H\) with exponent \(\kappa\) and constant \(B\), relative to
\(h^{\circ}\). Let \((X_1,Y_1),\ldots,(X_m,Y_m)\) be iid with law \(P\),
\(m\geq1\), and suppose \(\mathcal T_H\) is self-normalized
symmetrization-\allowbreak measurable for \(P\) at sample size \(m\); a
pointwise countable core suffices
(\cref{thm:p15-w02b-p15-065-ternary-core}). For \(0<\delta<1\) put
\[
 L=\ln\frac{4\Pi_H(2m)}{\delta}.
\]
Outside a set of outer probability at most \(\delta\), every \(h\in H\)
satisfies
\[
 R(h)-R_H^{*}\leq\widehat R_m(h)-\widehat R_m(h^{\circ})
 +\sqrt{\frac{160B\,(R(h)-R_H^{*})^{\kappa}L}{m}}+\frac{80L}{m}.
\]
Outside the same set:
\begin{enumerate}
\item every \(h\in H\) with \(\widehat R_m(h)\leq\widehat R_m(h^{\circ})\),
      in particular every exact empirical minimizer, satisfies
      \[
       R(h)-R_H^{*}\leq\max\Bigl\{\Bigl(\frac{640BL}{m}\Bigr)^{1/(2-\kappa)},\
       \frac{160L}{m}\Bigr\};
      \]
\item every \(\vartheta\)-approximate empirical minimizer,
      \(\vartheta\geq0\), satisfies
      \[
       R(h)-R_H^{*}\leq\max\Bigl\{3\vartheta,\
       \Bigl(\frac{1440BL}{m}\Bigr)^{1/(2-\kappa)},\ \frac{240L}{m}\Bigr\}.
      \]
\end{enumerate}
If the VC dimension \(d\) of the positive prediction-event class satisfies
\(1\leq d\leq2m\), then \(L\leq d\ln(2em/d)+\ln(4/\delta)\) by
\cref{thm:p15-vcl-sauer}. The selection among minimizers may be arbitrary
and need not be measurable.
\end{hbtheorem}

\begin{hbcorollary}{thm:p15-w02b-p15-065-vc-noise-orders}{proved=annex:p15-w02b-p15-065-vc-noise-orders}{The three noise conditions for a class of finite VC dimension}
In the setting of \cref{thm:p15-w02b-p15-065-vc-noise-rate}, suppose \(H\)
contains a Bayes predictor \(f^{*}\), take \(h^{\circ}=f^{*}\), so that
\(R_H^{*}=R^{*}\), and let \(L\) be as there. Outside a set of outer
probability at most \(\delta\), every exact empirical minimizer
\(\widehat h\) satisfies the following.
\begin{enumerate}
\item Realizable by a member of \(H\): \(R(\widehat h)\leq640L/m\).
\item Bounded noise with margin \(\gamma\):
      \(R(\widehat h)-R^{*}\leq640L/(\gamma m)\).
\item Margin condition with exponent \(a\) and constant \(C\):
      \[
       R(\widehat h)-R^{*}\leq\max\Bigl\{\Bigl(\frac{1280C^{1/(1+a)}L}{m}\Bigr)^{(1+a)/(2+a)},\
       \frac{160L}{m}\Bigr\}.
      \]
\item No condition on the noise:
      \(R(\widehat h)-R^{*}\leq\max\{\sqrt{640L/m},\,160L/m\}\).
\end{enumerate}
With \(1\leq d\leq2m\) and \(L\leq d\ln(2em/d)+\ln(4/\delta)\), the
bound on the excess risk is thus of order \(L/m\) in items~1 and~2, of
order \((L/m)^{(1+a)/(2+a)}\) in item~3, and of order \((L/m)^{1/2}\) in
item~4, where \(L\) grows like \(d\ln m\).
\end{hbcorollary}

\begin{hbexample}{def:p15-w02b-p15-065-vc-noise-example}{All thresholds with bounded noise}
Keep the law of \cref{def:p15-w02b-p15-065-noise-example}, and replace the
\(101\) thresholds by all of them:
\(H=\{h_s:s\in[0,1]\}\) with \(h_s(x)=\mathbf1\{x\geq s\}\).
\Cref{thm:p15-w02b-p15-065-vc-noise-example} proves the following. The
positive prediction-event class has a pointwise countable core and
\(\Pi_H(t)\leq t+1\). The member \(h_{3/10}\) is a Bayes predictor, and the
variance condition holds on \(H\) relative to it with \(\kappa=1\) and
\(B=5/4\). With \(\delta=1/20\) and \(m=1550000\), item~2 of
\cref{thm:p15-w02b-p15-065-vc-noise-orders} gives
\(R(\widehat h)-R^{*}\leq0.01\), that is \(R(\widehat h)\leq0.11\), for
every exact empirical minimizer outside a set of outer probability at most
\(0.05\). At the same \(m\) and \(\delta\), the bound of item~4, which
ignores the noise structure, exceeds \(0.089\).
\end{hbexample}

\begin{hbtheorem}{thm:p15-w02b-p15-065-vc-noise-example}{proved=annex:p15-w02b-p15-065-vc-noise-example}{The calculations of the example with all thresholds}
Every assertion of \cref{def:p15-w02b-p15-065-vc-noise-example} holds.
\end{hbtheorem}

The sample size of this example is about two hundred times that of the
finite example, for the same accuracy and confidence. The difference comes
from the constants of the ghost-sample argument and from the factor
\(\ln(4\Pi_H(2m)/\delta)\), which grows with \(m\). The theorem gives a
sufficient size, not the smallest one, and its factor \(d\ln(2em/d)\) is
not shown here to be necessary.

\begin{hbwarning}{def:p15-w02b-p15-065-noise-warning}{What a noise condition requires}
The bounded-noise and margin conditions are statements about the
conditional label probability \(\eta\), which the sample does not reveal;
they are assumptions about the law, not conclusions from data. The sizes
of items~1 to~3 of \cref{thm:p15-w02b-p15-065-noise-sizes} also need a
Bayes predictor in the finite class. If the class misses it, the transfer
of \cref{thm:p15-w02b-p15-065-noise-relations} does not apply, and the
variance condition must be checked on the class relative to its best
member. If the law puts positive probability on \(\{\eta=1/2\}\), item~7
of that lemma excludes every exponent \(\kappa>0\) on all measurable
predictors. For a class of finite VC dimension the rates of
\cref{thm:p15-w02b-p15-065-vc-noise-rate} also need the measurability
hypothesis of \cref{def:p15-w02b-p15-065-ternary}; without it, or a
countable core, the theorem draws no conclusion.
\end{hbwarning}

\section{Procedure, limits, relationships and sources}
\label{sec:p15-w02b-p15-065-procedure}

\begin{hbmethod}{def:p15-w02b-p15-065-method}{From a class to a guarantee, or to its impossibility}
\begin{enumerate}
\item Fix the class and the loss before sampling. For a loss with values in
      \([a,b]\), form the level-set class
      (\cref{def:p15-w02b-p15-065-level-class}); for binary loss it is the
      loss-event class.
\item Establish measurability: a countable core
      (\cref{thm:p15-w02b-p15-065-core-sufficient}), or a direct proof
      that \(D_m\) and \(K_m\) are measurable for the completed law, as in
      \cref{thm:p15-w02b-p15-065-singletons}. If neither is available, no
      uniform conclusion is drawn.
\item Bound the growth by Sauer's lemma (\cref{thm:p15-vcl-sauer}) and
      apply \cref{thm:p15-w02b-p15-065-outer-uniform} or
      \cref{thm:p15-w02b-p15-065-real-uniform}.
\item Decide the regime. Use the zero-risk bound only if a member of risk
      zero exists; a positive noise floor rules this out
      (\cref{thm:p15-w02b-p15-065-noise}). For a class containing a
      Bayes predictor, a bounded-noise or margin assumption on the law
      gives the intermediate sizes of
      \cref{thm:p15-w02b-p15-065-noise-sizes} for a finite class and the
      intermediate rates of \cref{thm:p15-w02b-p15-065-vc-noise-orders}
      for a class of finite VC dimension satisfying the measurability
      hypothesis of \cref{def:p15-w02b-p15-065-ternary}; state the
      assumption, since the sample does not establish it.
\item Report the conclusion in outer probability, over all members and all
      laws for which step~2 holds. For a validation estimate, require a
      jointly measurable learner and a validation sample not used in
      selection.
\item To argue that fewer samples cannot suffice for all laws, exhibit a
      shattered set and apply the lower bounds of
      \cref{sec:p15-w02b-p15-065-lower}.
\end{enumerate}
\end{hbmethod}

\emph{Limitations.} All results assume an independent sample from one fixed
law and a class fixed before sampling. The repeated-input and changed-law
countercases of \cref{thm:p15-occam-boundaries} remain outside every bound
above. When the law at use differs from the sampling law, a risk bound
proved here transfers only under a stated control of the difference, as in
\cref{thm:p15-w02b-p15-066-tv-transport,thm:p15-w02b-p15-066-density-transport}.
The constants are sufficient, not sharp. The lower bounds are
minimax over finitely supported laws; they say nothing about a fixed law.
The intermediate noise rates of
\cref{sec:p15-w02b-p15-065-noise-conditions} rest on an assumption about
the conditional label probability that no sample verifies. For a class of
finite VC dimension they carry the factor \(d\ln(2em/d)\) of the growth
bound and constants eighty times those of the finite-class theorem; whether the logarithmic factor can be removed is not settled
here.

\emph{Relationships.} The countable-core route of
\cref{sec:p15-vc-dimension-and-uniform-learning} is the special case of
\cref{thm:p15-w02b-p15-065-outer-uniform} in which all events are
measurable. The finite-class bounds of \cref{thm:p15-conc-finite-learning}
and the code-weighted bounds of \cref{sec:p15-occam-learning} use cardinality
or description length instead of traces. The decomposition of excess risk
into estimation and approximation terms is
\cref{thm:p15-w02b-p15-065-noise}, item~3, specialized from
\cref{thm:p15-w02b-p15-064-decomposition}. The noise conditions of
\cref{sec:p15-w02b-p15-065-noise-conditions} sit between the realizable
consistency bound and the agnostic bound: the variance condition with
\(\kappa=1\) recovers the order \(1/\varepsilon\) of
\cref{thm:p15-w02b-p15-065-real-learning}, item~3, and with \(\kappa=0\)
the order \(1/\varepsilon^2\) of \cref{thm:p15-conc-finite-learning}.
The rate proof replaces Hoeffding's range-only bound by the Bernstein form
of \cref{thm:p15-variance-tail-chain}, which uses the variance. For a
class of finite VC dimension the same role is played by
\cref{thm:p15-w02b-p15-065-relative}, which runs the ghost-sample and swap
argument of \cref{thm:p15-w02b-p15-065-outer-uniform} with the swap bound
normalized by the number of pooled points where a member is nonzero. The
optimality of the divergence bound within the exponential method, and its
comparison with the Chebyshev and normal benchmarks, are
\cref{thm:p15-w02b-p15-063-exponential-optimality,thm:p15-w02b-p15-063-chebyshev-comparison,thm:p15-w02b-p15-063-normal-benchmark}.

\emph{Sources.} The uniform-convergence question for frequencies of events
is posed by \cite{VapnikChervonenkis1971Uniform}; that paper's own
estimates and its necessary-and-sufficient entropy condition are not
reproduced here. The ghost-sample comparison is stated in
\cite[Section~5.2, Lemma~4, p.~17]{LuxburgScholkopf2008Author} without
proof or measurability conditions. The outer-probability regime, the
level-set reduction, the noise conditions and their rates, the examples and
the lower bounds above are proved in
full in the annex, and no source is credited with their constants.
